\chapter{Correspondance détaillée des catégories}\label{ann:correspondance}
Les statuts « fiable », « partielle » et « absente » caractérisent la correspondance à l'échelle de la division mensuelle. Le statut « fiable » reprend la catégorie \texttt{proche} du CSV d'audit. L'exclusion du groupe 14 répond en plus à sa nature économique. Les prix associés aux groupes partiels servent uniquement au scénario élargi.
\begin{table}[!ht]
\centering\small
\setlength{\tabcolsep}{4pt}
\caption{Audit des correspondances ENNVM--IPC (groupes 01 à 07)}\label{tab:correspondance}
\begin{tabular}{@{}l>{\raggedright\arraybackslash}p{.235\textwidth}>{\raggedright\arraybackslash}p{.21\textwidth}l>{\raggedright\arraybackslash}p{.35\textwidth}@{}}
\hline
N° & Groupe ENNVM & Division IPC & Statut & Justification \\
\hline
01 & Aliments et boissons non alcoolisées & Produits alimentaires et boissons non alcoolisées & Fiable & DAMPG01 et IPC 011--012 : aliments (dont céréales, viandes et fruits) et boissons sans alcool. La division 01 est plus ciblée que l'agrégat Alimentation. \\
\hline
02 & Boissons alcoolisées, tabac et stupéfiants & Boissons alcoolisées, tabac et stupéfiants & Partielle & DAMPG02 publie une dépense de drogue (23), distincte de l'alcool et du tabac. Le détail IPC consulté montre 021 et 022, sans série identifiable pour les stupéfiants. \\
\hline
03 & Habillement et chaussures & Articles d'habillement et chaussures & Fiable & DAMPG03 et IPC 031--032 : vêtements, chaussures et services associés comme le nettoyage, la réparation ou la cordonnerie. \\
\hline
04 & Logement, eau, électricité, gaz et autres combustibles & Logement, eau, gaz, électricité et autres combustibles & Partielle & DAMPG04 comprend les loyers imputés (42). L'IPC documente les loyers payés par les locataires (041), avec recoupement sur l'eau, l'énergie et l'entretien. \\
\hline
05 & Ameublement, équipement ménager et entretien courant de la maison & Meubles, articles de ménage et entretien courant du foyer & Fiable & DAMPG05 et IPC 051--056 : ameublement, textiles, appareils, vaisselle, outils et entretien courant. \\
\hline
06 & Santé & Santé & Partielle & Médicaments, soins ambulatoires et hospitaliers se recoupent. DAMPG06 inclut aussi les groupes 64 et 65 (dont médecine traditionnelle), sans équivalent identifiable pour l'ensemble de ces dépenses dans le détail IPC consulté. \\
\hline
07 & Transport & Transports & Partielle & Véhicules, utilisation et transport de voyageurs se recoupent. DAMPG07 inclut aussi le fret et les services postaux (74), ces derniers figurant dans l'IPC 081. \\
\hline
\end{tabular}
\sourcetable{Fichier de correspondance du projet, groupes DAMPG de l'ENNVM et détail des classes de l'IPC.}
\end{table}
\clearpage
\begin{table}[!ht]
\centering\small
\setlength{\tabcolsep}{4pt}
\caption{Audit des correspondances ENNVM--IPC (groupes 08 à 14)}
\begin{tabular}{@{}l>{\raggedright\arraybackslash}p{.235\textwidth}>{\raggedright\arraybackslash}p{.21\textwidth}l>{\raggedright\arraybackslash}p{.35\textwidth}@{}}
\hline
N° & Groupe ENNVM & Division IPC & Statut & Justification \\
\hline
08 & Information et communication & Communications & Partielle & Recoupement sur la téléphonie. DAMPG08 comprend aussi le matériel informatique et audiovisuel que le détail IPC classe sous Loisirs et culture (091). \\
\hline
09 & Loisirs, sport et culture & Loisirs et culture & Partielle & Le sport, la culture et la papeterie se recoupent. L'IPC 091 inclut le matériel audiovisuel et informatique qui relève de DAMPG08 dans l'enquête. \\
\hline
10 & Enseignement & Enseignement & Fiable & DAMPG10 et IPC 101--105 : préscolaire et primaire, secondaire, postsecondaire non supérieur, supérieur et enseignement non défini par niveau. \\
\hline
11 & Services de restauration et d'hébergement (S) & Restaurants et hôtels & Partielle & Restauration, cantines et hébergement figurent des deux côtés. DAMPG11 détaille aussi les repas livrés et divers services dont la couverture précise reste incertaine. \\
\hline
12 & Assurances et services financiers & Aucune & Absente & Aucune série mensuelle autonome. Le classeur IPC annuel de détail situe les assurances transports et l'intermédiation financière dans Biens et services divers. \\
\hline
13 & Soins personnels, protection sociale et biens et services divers & Biens et services divers & Partielle & Recoupement sur les soins et effets personnels. DAMPG13 inclut la protection sociale et la division IPC recoupe aussi le groupe 12. \\
\hline
14 & Paiements fiscaux, transferts et remboursement des prêts (Non liés à l'exercice d'activité professionnelle) & Aucune & Absente & Poste distinct exclu des deux scénarios. Les paiements fiscaux, transferts et remboursements ne sont pas nécessairement des achats de biens et services. \\
\hline
\end{tabular}
\sourcetable{Fichier de correspondance du projet, groupes DAMPG de l'ENNVM et détail des classes de l'IPC.}
\end{table}


\chapter{Inventaire et contrôles de reproductibilité}\label{ann:trace}
\begin{description}\raggedright
\item[Données de prix.] Trois classeurs mensuels du HCP, base 100 en 2017, puis \path{ipc_maroc_mensuel.csv} (118 mois, 15 séries et une date). Le classeur annuel de détail des services divers documente certaines classes.
\item[Données de dépenses.] Classeur de tabulation ENNVM 2022--2023.
\par \path{depenses_menages_2022_2023.csv} (252 couples profil--poste) et \path{profils_budgetaires_2022_2023.csv} (18 profils).
\item[Correspondance.] Audit des 14 groupes dans \path{correspondance_depenses_ipc.csv}. Le notebook 03 contient le rapport de préparation, les couvertures et la définition de \texttt{simuler\_panier}.
\item[Analyse de l'IPC.] Le notebook 02 fournit les statistiques et quatre graphiques préfixés \texttt{02\_}. Le notebook 01 conserve la trame exploratoire, sans résultats exécutés.
\item[Simulations.] \path{04_cout_panier_historique.csv} contient 4\,248 lignes (118 mois, 18 profils, deux périmètres). \path{04_chocs_hypothetiques.csv} contient 72 lignes (deux chocs pour chaque couple profil--périmètre). \path{04_ponderations_sous_paniers.csv} contient 288 poids (18 profils, quatre postes prudents et douze élargis). Deux graphiques du notebook 04 complètent ces tableaux.
\item[Scénarios de ressources.] Le notebook 05 croise les coûts du notebook 04 avec quatre taux nominaux hypothétiques. Les fichiers \path{05_pouvoir_achat_historique.csv} (16\,992 lignes) et \path{05_pouvoir_achat_chocs.csv} (288 lignes) conservent la couverture et le seuil de maintien.
\item[Robustesse.] Le notebook 06 compare les dates de référence et les variantes de pondération alimentaire. Il enregistre ses tableaux et figures dans \path{outputs/robustesse/}.
\end{description}

Les contrôles du notebook 04 portent sur les clés, la complétude du calendrier, les valeurs finies, la conservation des poids sources, la normalisation à 100 en septembre 2022, et la concordance entre contribution directe et contribution obtenue via la couverture. Le choc uniforme produit exactement 5\,\% à la précision du calcul. Le supplément alimentaire marginal est identique entre périmètres pour un profil donné. Les taux historiques ont été confrontés à \texttt{simuler\_panier}. Les montants affichés dans le rapport sont arrondis après calcul.

\clearpage
\chapter{Tableaux complets des chocs hypothétiques}\label{ann:chocs}
Les colonnes de contribution donnent des points du budget initial. Les montants illustrent un même budget de 5\,000 MAD pour chaque profil. « Combiné » correspond à +5\,\% sur tous les postes couverts, puis +10\,\% supplémentaires sur le groupe 01.
\begin{table}[!htbp]
\centering\small
\caption{Contributions au budget initial, périmètre prudent}\label{tab:contributions-prudent}
\begin{tabular}{llrrr}
\hline
Milieu & Classe & Uniforme (pts) & Combiné (pts) & Surcroît aliment. (MAD) \\
\hline
National & C1 & 2{,}60 & 7{,}42 & 241{,}50 \\
National & C2 & 2{,}44 & 6{,}91 & 223{,}65 \\
National & C3 & 2{,}31 & 6{,}47 & 207{,}90 \\
National & C4 & 2{,}15 & 5{,}93 & 189{,}00 \\
National & C5 & 1{,}80 & 4{,}61 & 140{,}18 \\
National & Ensemble & 2{,}07 & 5{,}58 & 175{,}35 \\
Urbain & C1 & 2{,}25 & 6{,}36 & 205{,}28 \\
Urbain & C2 & 2{,}17 & 6{,}05 & 193{,}72 \\
Urbain & C3 & 2{,}12 & 5{,}81 & 184{,}80 \\
Urbain & C4 & 2{,}02 & 5{,}44 & 171{,}15 \\
Urbain & C5 & 1{,}73 & 4{,}31 & 129{,}15 \\
Urbain & Ensemble & 1{,}93 & 5{,}09 & 158{,}03 \\
Rural & C1 & 2{,}79 & 8{,}07 & 263{,}55 \\
Rural & C2 & 2{,}75 & 7{,}89 & 256{,}73 \\
Rural & C3 & 2{,}70 & 7{,}74 & 252{,}00 \\
Rural & C4 & 2{,}63 & 7{,}50 & 243{,}60 \\
Rural & C5 & 2{,}33 & 6{,}59 & 213{,}15 \\
Rural & Ensemble & 2{,}55 & 7{,}28 & 236{,}25 \\
\hline
\end{tabular}
\sourcetable{Simulations hypothétiques sur un budget initial de 5\,000 MAD.}
\end{table}

\begin{table}[!htbp]
\centering\small
\caption{Contributions au budget initial, périmètre élargi}\label{tab:contributions-elargi}
\begin{tabular}{llrrr}
\hline
Milieu & Classe & Uniforme (pts) & Combiné (pts) & Surcroît aliment. (MAD) \\
\hline
National & C1 & 4{,}90 & 9{,}73 & 241{,}50 \\
National & C2 & 4{,}86 & 9{,}33 & 223{,}65 \\
National & C3 & 4{,}79 & 8{,}95 & 207{,}90 \\
National & C4 & 4{,}72 & 8{,}51 & 189{,}00 \\
National & C5 & 4{,}50 & 7{,}31 & 140{,}18 \\
National & Ensemble & 4{,}66 & 8{,}17 & 175{,}35 \\
Urbain & C1 & 4{,}88 & 8{,}99 & 205{,}28 \\
Urbain & C2 & 4{,}81 & 8{,}68 & 193{,}72 \\
Urbain & C3 & 4{,}75 & 8{,}45 & 184{,}80 \\
Urbain & C4 & 4{,}67 & 8{,}09 & 171{,}15 \\
Urbain & C5 & 4{,}47 & 7{,}06 & 129{,}15 \\
Urbain & Ensemble & 4{,}62 & 7{,}78 & 158{,}03 \\
Rural & C1 & 4{,}92 & 10{,}19 & 263{,}55 \\
Rural & C2 & 4{,}89 & 10{,}02 & 256{,}73 \\
Rural & C3 & 4{,}85 & 9{,}89 & 252{,}00 \\
Rural & C4 & 4{,}79 & 9{,}66 & 243{,}60 \\
Rural & C5 & 4{,}63 & 8{,}90 & 213{,}15 \\
Rural & Ensemble & 4{,}76 & 9{,}48 & 236{,}25 \\
\hline
\end{tabular}
\sourcetable{Simulations hypothétiques sur un budget initial de 5\,000 MAD.}
\end{table}


\clearpage
\chapter{Pouvoir d'achat après le choc combiné}\label{ann:pa-chocs}
Le tableau de chaque périmètre présente les quatre hypothèses nominales et le seuil de maintien. Le groupe alimentaire subit une hausse totale de 15,5\,\% et les autres postes couverts de 5\,\%.
\begin{table}[!htbp]
\centering\small
\caption{Après le choc hypothétique combiné, périmètre prudent}\label{tab:pa-c-prudent}
\begin{tabular}{llrrrrrr}
\hline
Milieu & Classe & Couvert & Seuil & 0\,\% & +3\,\% & +5\,\% & +10\,\% \\
\hline
National & C1 & 51{,}9 & 14{,}31 & -12{,}52 & -9{,}89 & -8{,}14 & -3{,}77 \\
National & C2 & 48{,}8 & 14{,}17 & -12{,}41 & -9{,}78 & -8{,}03 & -3{,}65 \\
National & C3 & 46{,}2 & 14{,}00 & -12{,}28 & -9{,}65 & -7{,}89 & -3{,}51 \\
National & C4 & 43{,}0 & 13{,}79 & -12{,}12 & -9{,}48 & -7{,}73 & -3{,}33 \\
National & C5 & 36{,}1 & 12{,}77 & -11{,}32 & -8{,}66 & -6{,}89 & -2{,}45 \\
National & Ensemble & 41{,}4 & 13{,}47 & -11{,}87 & -9{,}23 & -7{,}47 & -3{,}06 \\
\noalign{\vskip 2pt}
Urbain & C1 & 45{,}0 & 14{,}12 & -12{,}38 & -9{,}75 & -7{,}99 & -3{,}61 \\
Urbain & C2 & 43{,}5 & 13{,}91 & -12{,}21 & -9{,}58 & -7{,}82 & -3{,}43 \\
Urbain & C3 & 42{,}3 & 13{,}74 & -12{,}08 & -9{,}44 & -7{,}68 & -3{,}29 \\
Urbain & C4 & 40{,}4 & 13{,}47 & -11{,}87 & -9{,}23 & -7{,}47 & -3{,}06 \\
Urbain & C5 & 34{,}6 & 12{,}47 & -11{,}08 & -8{,}42 & -6{,}64 & -2{,}19 \\
Urbain & Ensemble & 38{,}6 & 13{,}19 & -11{,}65 & -9{,}00 & -7{,}23 & -2{,}82 \\
\noalign{\vskip 2pt}
Rural & C1 & 55{,}9 & 14{,}43 & -12{,}61 & -9{,}99 & -8{,}24 & -3{,}87 \\
Rural & C2 & 55{,}1 & 14{,}32 & -12{,}53 & -9{,}90 & -8{,}15 & -3{,}78 \\
Rural & C3 & 54{,}0 & 14{,}33 & -12{,}54 & -9{,}91 & -8{,}16 & -3{,}79 \\
Rural & C4 & 52{,}6 & 14{,}26 & -12{,}48 & -9{,}86 & -8{,}11 & -3{,}73 \\
Rural & C5 & 46{,}5 & 14{,}17 & -12{,}41 & -9{,}78 & -8{,}03 & -3{,}65 \\
Rural & Ensemble & 51{,}0 & 14{,}26 & -12{,}48 & -9{,}86 & -8{,}11 & -3{,}73 \\
\hline
\end{tabular}
\sourcetable{Notebook 05. Couvert : part du budget initial en pourcentage. Seuil : croissance nominale de maintien en pourcentage. Les quatre dernières colonnes donnent la variation des unités accessibles (en \%).}
\end{table}

\begin{table}[!htbp]
\centering\small
\caption{Après le choc hypothétique combiné, périmètre élargi}\label{tab:pa-c-elargi}
\begin{tabular}{llrrrrrr}
\hline
Milieu & Classe & Couvert & Seuil & 0\,\% & +3\,\% & +5\,\% & +10\,\% \\
\hline
National & C1 & 98{,}0 & 9{,}93 & -9{,}03 & -6{,}30 & -4{,}48 & 0{,}06 \\
National & C2 & 97{,}1 & 9{,}61 & -8{,}76 & -6{,}03 & -4{,}20 & 0{,}36 \\
National & C3 & 95{,}9 & 9{,}34 & -8{,}54 & -5{,}79 & -3{,}97 & 0{,}61 \\
National & C4 & 94{,}5 & 9{,}00 & -8{,}26 & -5{,}50 & -3{,}67 & 0{,}92 \\
National & C5 & 90{,}1 & 8{,}11 & -7{,}50 & -4{,}73 & -2{,}88 & 1{,}75 \\
National & Ensemble & 93{,}2 & 8{,}76 & -8{,}06 & -5{,}30 & -3{,}46 & 1{,}14 \\
\noalign{\vskip 2pt}
Urbain & C1 & 97{,}7 & 9{,}20 & -8{,}43 & -5{,}68 & -3{,}85 & 0{,}73 \\
Urbain & C2 & 96{,}2 & 9{,}03 & -8{,}28 & -5{,}53 & -3{,}69 & 0{,}89 \\
Urbain & C3 & 95{,}1 & 8{,}89 & -8{,}16 & -5{,}41 & -3{,}57 & 1{,}02 \\
Urbain & C4 & 93{,}4 & 8{,}66 & -7{,}97 & -5{,}21 & -3{,}37 & 1{,}23 \\
Urbain & C5 & 89{,}5 & 7{,}89 & -7{,}31 & -4{,}53 & -2{,}68 & 1{,}96 \\
Urbain & Ensemble & 92{,}4 & 8{,}42 & -7{,}77 & -5{,}00 & -3{,}15 & 1{,}46 \\
\noalign{\vskip 2pt}
Rural & C1 & 98{,}3 & 10{,}36 & -9{,}39 & -6{,}67 & -4{,}86 & -0{,}33 \\
Rural & C2 & 97{,}8 & 10{,}25 & -9{,}30 & -6{,}58 & -4{,}76 & -0{,}23 \\
Rural & C3 & 97{,}0 & 10{,}20 & -9{,}25 & -6{,}53 & -4{,}72 & -0{,}18 \\
Rural & C4 & 95{,}8 & 10{,}09 & -9{,}16 & -6{,}44 & -4{,}62 & -0{,}08 \\
Rural & C5 & 92{,}7 & 9{,}60 & -8{,}76 & -6{,}02 & -4{,}20 & 0{,}37 \\
Rural & Ensemble & 95{,}2 & 9{,}96 & -9{,}06 & -6{,}33 & -4{,}51 & 0{,}03 \\
\hline
\end{tabular}
\sourcetable{Notebook 05. Couvert : part du budget initial en pourcentage. Seuil : croissance nominale de maintien en pourcentage. Les quatre dernières colonnes donnent la variation des unités accessibles (en \%).}
\end{table}

